\documentclass[10pt,a4paper]{amsart}
\usepackage[margin=19mm]{geometry}
\usepackage{amsmath,amssymb,mathtools}
\usepackage[scr=rsfs,cal=euler]{mathalfa}
\usepackage{xfrac}
\usepackage[unicode,hidelinks,bookmarks=false]{hyperref}
\newcommand{\ie}{i.e.\ }
\newcommand{\cf}{cf.\ }
\newcommand{\resp}{respectively\ }
\newcommand{\Subsection}{\subsection}
\providecommand{\nobreakdash}{\nobreak\hbox{-}}
\setlength{\emergencystretch}{2em}
\multlinegap0pt
\usepackage{bm}
\usepackage{bbm}
\usepackage{dsfont}
\usepackage{xspace}
\usepackage{cancel}
\usepackage{mathtools}
\usepackage{amsthm}
\makeatletter
\@ifundefined{theorem}{\newtheorem{theorem}{Theorem}}{}
\@ifundefined{lemma}{\newtheorem{lemma}{Lemma}}{}
\makeatother
\title[Likelihood-based Spacings Goodness-of-Fit Statistics for Uni]{Likelihood-based Spacings Goodness-of-Fit Statistics for Univariate Shape-constrained Densities}
\author{Kwun Chuen Gary Chan, Hok Kan Ling, Chuan-Fa Tang, Sheung Chi Phillip Yam}
\date{Source: arXiv:2211.13272v3 (CC BY 4.0)}
\begin{document}
\maketitle

\noindent\emph{Source and licence.} Likelihood-based Spacings Goodness-of-Fit Statistics for Univariate Shape-constrained Densities. Version arXiv:2211.13272v3 (CC BY 4.0), distributed under the Creative Commons Attribution 4.0 International licence, as recorded in the arXiv licence metadata for this exact version and shown on the abstract page https://arxiv.org/abs/2211.13272v3. Full rights notice: licenses/arxiv-2211.13272-CC-BY-4.0.txt.\par\smallskip

\noindent\textit{Editorial scope.} Counted unit: the Lemma whose source label is \texttt{lemma:Rn\_terms} in the article (source0.tex lines 1569--1576, source label \texttt{lemma:Rn\_terms}), delivered together with the article's own complete original proof of it (source0.tex lines 1578--1706, one \texttt{proof} environment of 129 lines). No mathematics is written, repaired or re-derived: statement and proof are the article's own text line for line. Required context retained uncounted: lemma:dist\_same (lines 1290--1301). Every definition, display and label of those lines is retained unchanged. Omitted from this package and not redistributed: 1--1289, 1302--1568, 1577--1577, 1707--3447. Nothing inside a retained excerpt is omitted or altered: every retained body is delivered exactly as the article's lines are written, so no omission receipt of any kind accompanies this package. Citations: the retained excerpts carry no citation command at all, so no bibliography entry of the article is transcribed and no reference is redistributed. Reference adaptation: none was needed and none was made; every reference inside the delivered unit resolves inside it, so no unresolved reference or citation is produced. Printed slips kept literal: no printed word, symbol, hypothesis or number of the counted statement or of its proof was altered. Layout and class: the article's own class is \texttt{article} with packages including geometry, lineno, lipsum, amsmath, amssymb, bm, natbib, amsfonts, xr, enumerate, graphicx, amsthm, xcolor, csquotes; the wrapper is the lane's pinned \texttt{amsart} editorial wrapper at 10pt on A4, which restates the theorem-like environments the retained text needs and adds the article's own macro definitions that the retained text needs (none), with extra wrapper packages (bm, bbm, dsfont, xspace, cancel, mathtools). The delivered numbering is the unit's own. Assets: no retained span contains a figure environment, a graphics-inclusion command, a PDF-page inclusion, a table or a TikZ picture; no image file is copied into this package, the package declares no asset dependency and no image file is redistributed with it. Licence: version arXiv:2211.13272v3 of this article is distributed under the Creative Commons Attribution 4.0 International licence, as recorded in the arXiv licence metadata for this exact version and shown on the abstract page https://arxiv.org/abs/2211.13272v3. A full rights notice is delivered at licenses/arxiv-2211.13272-CC-BY-4.0.txt. Characters: every retained excerpt is pure ASCII, so the delivered engine, XeLaTeX with its default Latin Modern text font, reports no missing character.
\begin{lemma}\label{lemma:dist_same}
	We have
	\begin{equation*}
		M_n - \frac{n-1}{n}\log \nu \stackrel{d}{=} M_{1n} + M_{2n},
	\end{equation*}
	where
	\begin{align*}
		M_{1n} &:= \frac{1}{n} \sum_{j=0,\ldots, \frac{n-1}{\nu}-1} \left( \tilde{E}_j - \nu - \nu \log \tilde{E}_j \right),\\
		M_{2n} &:= \frac{1}{n-1} (E_1+E_{n+1})- \frac{Y_n}{n} - (1-n^{-1})\frac{Y_n Y^*_n}{1+Y^*_n},
	\end{align*}
	and $Y^*_n$ lies between $0$ and $Y_n$.
\end{lemma}

\begin{lemma}\label{lemma:Rn_terms}
Suppose one of Conditions (C) (i) or (ii) holds.
Then, we have
\begin{equation*}
    |R_n| \precsim \frac{\nu}{n}(C + |\log f_0(Z_1)| + |\log f_0(Z_n)|) +O_p\left(\frac{\log n}{n}\right),
\end{equation*}
for some constant $C$ that depends on $K_1,K_2, L$ if Condition (ii) holds and is a universal constant if Condition (i) holds.
\end{lemma}

\begin{proof}[Proof of Lemma \ref{lemma:Rn_terms}]
	Write
\begin{equation*}
	R_n = R_{1n} + R_{2n},
\end{equation*}
where
\begin{align*}
R_{1n} &:=	- \frac{1}{n} \sum_{j=0,\ldots,\frac{n-1}{\nu} - 1} 
	\sum^\nu_{l=1} \log \frac{f_0(Z_{j\nu+l+1})(Z_{(j+1)\nu+1} - Z_{j\nu+1})}{F_0(Z_{(j+1)\nu+1}) - F_0(Z_{j\nu+1})},\\
	R_{2n} &:= -\frac{1}{n}\log \frac{f_0(Z_1)}{f^H_{n,\nu}(Z_1)}.
\end{align*}

%
For $R_{1n}$, by the mean value theorem, there exists $Z^*_j$ lying between $Z_{j\nu+1}$ and $Z_{(j+1)\nu+1}$ such that $F_0(Z_{(j+1)\nu+1}) - F_0(Z_{j\nu+1}) = f_0(Z^*_j)(Z_{(j+1)\nu+1} - Z_{j\nu+1})$ and so
	\begin{equation*}
R_{1n}  = -\frac{1}{n}\sum_{j=0,\ldots, \frac{n-1}{\nu}-1} \sum^\nu_{l=1} \log \frac{f_0(Z_{j\nu+l+1})}{f_0(Z^*_j)}.
	\end{equation*}
	%
 If Condition (i) holds and $f_0$ is monotone increasing, then
\begin{align}
    |n R_{1n}| & \leq \left|  \sum_{j=0,\ldots, \frac{n-1}{\nu}-1} \sum^\nu_{l=1} \log \frac{f_0(Z_{j\nu+1})}{f_0(Z_{(j+1)\nu+1})}\right| =\nu \left| \log f_0(Z_1) - \log f_0(Z_n)\right|\nonumber \\
    & \leq \nu (|\log f_0(Z_1)| + |\log f_0(Z_n)|). \label{eqnR_1n1}
\end{align}
We have the same inequality when $f_0$ is monotone decreasing.
If Condition (ii) holds, 
let 
\begin{equation*}
R^{(m)}_{1n} := \sum_{j=0,\ldots, \frac{n-1}{\nu}-1} \sum^\nu_{l=1} I^{(m)}_{n,jl} \log \frac{f_0(Z_{j\nu+l+1})}{f_0(Z^*_j)},
\end{equation*}
where 
\begin{align*}
    I^{(1)}_{n,jl} &:= \mathbbm{1}(Z_{j\nu+l+1}, Z^*_j \in [K_1,K_2]), \quad 
    I^{(2)}_{n,jl} := \mathbbm{1}( Z^*_j < K_2 < Z_{j\nu+l+1}),\\
    I^{(3)}_{n,jl} &:= \mathbbm{1}(  Z_{j\nu+l+1} < K_2 < Z^*_j), \quad 
    I^{(4)}_{n,jl}:=\mathbbm{1}( Z^*_j < K_1 < Z_{j\nu+l+1}),\\
    I^{(5)}_{n,jl} &:= \mathbbm{1}(  Z_{j\nu+l+1} < K_1 < Z^*_j), \quad 
    I^{(6)}_{n,jl} :=\mathbbm{1}(  Z_{j\nu+l+1}, Z^*_j > K_2),\\
    I^{(7)}_{n,jl} &:=\mathbbm{1}(  Z_{j\nu+l+1}, Z^*_j < K_1).
\end{align*}
Then,
\begin{equation*}
|nR_{1n}| \leq  \sum^7_{m=1}|R^{(m)}_{1n}|.
\end{equation*}
%
For $R^{(1)}_{1n}$, by the Lipschitz continuity of $\log f_0$ on $[K_1, K_2]$, 
\begin{align*}
    |R^{(1)}_{1n}|& \leq \sum_{j=0,\ldots, \frac{n-1}{\nu}-1} \sum^\nu_{l=1}  I^{(1)}_{n,jl} L \left| Z_{j\nu+l+1} - Z^*_j  \right|\\
    & \leq \sum_{j=0,\ldots, \frac{n-1}{\nu}-1} \sum^\nu_{l=1}  I^{(1)}_{n,jl} L(Z_{(j+1)\nu+1} - Z^*_j ) \\
    & \leq  \nu L(K_2 - K_1),
\end{align*}
where the second and third inequalities follow from the fact that $Z^*_j$ is between $Z_{j\nu+1}$ and $Z_{(j+1)\nu+1}$.

		
For $R^{(2)}_{1n}$, note that there will be at most one $j$, say, $\tilde{j}$, such that $\mathbbm{1}( Z^*_{\tilde{j}} < K_2 < Z_{\tilde{j}\nu+l+1})$ for some $l$. If $Z^*_{\tilde{j}} > K_1$, by telescoping and the triangle inequality,
\begin{align*}
    |R^{(2)}_{1n}|	 & \leq \sum^\nu_{l=1}\left( |\log f_0(Z_{\tilde{j} \nu+l+1}) - \log f_0(K_2)| + |\log f_0(K_2) - \log f_0(Z^*_j)|\right)\\
    &\leq \nu (|\log f_0(Z_n)| + |\log f_0(K_2)| + L(K_2 - K_1)),
\end{align*}
		where the last inequality follows from the monotonicity of $f_0$ beyond $K_2$ and the Lipschitz continuity of $\log f_0$ on $[K_1, K_2]$. 
		If $Z^*_{\tilde{j}} < K_1$, we have
		\begin{align*}
			|R^{(2)}_{1n}|	 & \leq \sum^\nu_{l=1}\bigg(  |\log f_0(Z_{\tilde{j} k+l+1}) - \log f_0(K_2)| + |\log f_0(K_2) - \log f_0(K_1)|\\
			& \quad  \quad  \quad  \quad + |\log f_0(K_1) - \log f_0(Z^*_j)|\bigg)\\
			&\leq \nu\bigg(|\log f_0(Z_n)| + |\log f_0(K_2)| + L(K_2-K_1) +  |\log f_0(K_1)| +  |\log f_0(Z_1)|\bigg).
		\end{align*}
		Thus, we always have
		\begin{align*}
			|R^{(2)}_{1n}| \leq \nu (C + |\log f_0(Z_1)|+|\log f_0(Z_n)|),
		\end{align*}
		where $C := L(K_2-K_1) + |\log f_0(K_2)| + |\log f_0(K_1)|$. It is straightforward to see the same bound holds for $|R^{(3)}_{1n}|, |R^{(4)}_{1n}|, |R^{(5)}_{1n}|$. The derivation of these bounds are similar to that for $|R^{(2)}_{1n}|$ and is therefore omitted.
		
		Now, we consider $R^{(6)}_n$. Let $j^*$ be the first $j$ such that $Z^*_j > K_2$. If $f_0$ is decreasing after $K_2$,
		\begin{align*}
			|R^{(6)}_{1n}|	&\leq \sum^\nu_{l=1} (\log f_0(K_2) - \log f_0(Z_{(j^*+1)\nu+1})) + \sum_{j=j^*+1,\ldots, \frac{n-1}{\nu}-1} \sum^\nu_{l=1}( \log f_0(Z_{j\nu+1}) - \log f_0(Z_{(j+1)\nu+1})) \\
			&\leq \nu(\log f_0(K_2) - \log f_0(Z_n) ).
		\end{align*}
		If $f_0$ is increasing after $K_2$, the bound becomes $\nu(\log f_0(Z_n) - \log f_0(K_2))$. In general, we have
		\begin{equation*}
			|R^{(6)}_{1n}| \leq \nu \left(|\log f_0(K_2)| + |\log f_0(Z_n)|\right).
		\end{equation*}
		Similarly, we also have
		\begin{equation*}
			|R^{(7)}_{1n}| \leq \nu\left(|\log f_0(K_1)| + |\log f_0(Z_1)|\right).
		\end{equation*}
		Combining the bounds for $|R^{(1)}_{1n}|,\ldots,|R^{(7)}_{1n}|$, we have
		\begin{equation}\label{eqnR_1n2}
			|nR_{1n}| \leq \nu(C_1 + 5|\log f_0(Z_1)| + 5|\log f_0(Z_n)|),
		\end{equation}
		where $C_1 := 5L(K_2-K_1) + 5|\log f_0(K_1)| + 5|\log f_0(K_2)|$. In view of (\ref{eqnR_1n1}) and (\ref{eqnR_1n2}), we have
  \begin{equation}\label{eq:R1n_order}
      |R_{1n}| \precsim \frac{v}{n}(C + |\log f_0(Z_1)| + |\log f_0(Z_n)|).
  \end{equation}

  For $R_{2n}$. Note that
\begin{align*}
	R_{2n} &= -\frac{1}{n} \log \frac{f_0(Z_1)(Z_{\nu+1} - Z_1)}{F_0(Z_{\nu+1})- F_0(Z_1)} - \frac{1}{n}\log \frac{F_0(Z_{\nu+1}) - F_0(Z_1))}{\nu/(n-1)}.
\end{align*}
By the mean value theorem, there exists $Z^*_1$ lying between $Z_{\nu+1}$ and $Z_1$ such that $F_0(Z_{\nu+1}) - F_0(Z_1) = f_0(Z^*_1)(Z_{\nu+1}-Z_1)$. Thus, using $(F_0(Z_1), F_0(Z_{\nu+1})) \stackrel{d}{=} (U_{(1)}, U_{(\nu+1)}) $ as in the proof of Lemma \ref{lemma:dist_same}, we have
\begin{equation}\label{eq:R2n1}
	R_{2n} \stackrel{d}{=}  -\frac{1}{n}\log \frac{f_0(Z_1)}{f_0(Z^*_1)} - \frac{1}{n} \log (U_{(\nu+1)} - U_{(1)}) + \frac{1}{n} \log \frac{n-1}{\nu}.
\end{equation}
Now, if Condition (i) holds, then 
\begin{equation}\label{eq:R2n2}
\left|	\log \frac{f_0(Z_1)}{f_0(Z^*_1)} \right| \leq \left|	\log \frac{f_0(Z_1)}{f_0(Z_n)} \right| \leq |\log f_0(Z_1)|+ |\log f_0(Z_n)|.
\end{equation}
If Condition (ii) holds, we have
\begin{align}
&	|\log f_0(Z^*_1)| \nonumber\\
 &= |(\log f_0(Z^*_1))\mathbbm{1}(Z^*_1 < K_1)  +
	(\log f_0(Z^*_1))\mathbbm{1}(K_1 < Z^*_1 < K_2) + (\log f_0(Z^*_1))\mathbbm{1}(Z^*_1 > K_2)| \nonumber \\
	& \leq (|\log f_0(Z_1)| + |\log f_0(K_1)|) + (L(K_2 - K_1) + |\log f_0(K_1)|) \nonumber \\
	& \quad \quad  + (|\log f_0(K_2)| + |\log f_0(Z_n)|)  \nonumber \\
	& =C_2 + |\log f_0(Z_1)| + |\log f_0(Z_n)|, \label{eq:R2n3}
\end{align}
where
\begin{equation*}
	C_2:= 2 |\log f_0(K_1)| + L(K_2-K_1) + |\log f_0(K_2)|.
\end{equation*}
Let $E_h$'s be i.i.d. standard exponential random variables the proof of Lemma \ref{lemma:dist_same}, then
\begin{align}
	|\log (U_{(\nu+1)} - U_{(1)})| & \leq |\log (U_{(2)} - U_{(1)})| \leq |\log E_2| + \left| \log \left( \sum^{n+1}_{h=1} E_h\right)\right|\nonumber\\
 &\leq |\log E_2| + \left| \log \left( \frac{1}{n+1}\sum^{n+1}_{h=1} E_h\right) \right|  + \log (n+1)= O_p(\log n).  \label{eq:R2n4}
\end{align}
In view of (\ref{eq:R2n1})-(\ref{eq:R2n4}), we have
\begin{equation}\label{eq:R2n_order}
\left|	R_{2n }\right| \leq \frac{2}{n} |\log f_0(Z_1)| + \frac{1}{n}|\log f_0(Z_n)| + O_p\left(\frac{\log n}{n}\right).
\end{equation}
The claim in the lemma follows in view of (\ref{eq:R1n_order}) and (\ref{eq:R2n_order}).
\end{proof}

\end{document}
